% SOURCE_STATEMENT: "In shear flows, why does transition to turbulence occur near the critical Reynolds number $R_{e,c}$ so often, and why is the resulting pattern's length scale so large?"
% Preparation and final self-audit: CONTRIBUTING.md prompts C and D; source/audit artifacts retained in the local physics_statements_2026_10_11 batch.
% !TeX program = lualatex
\documentclass[11pt]{article}
\usepackage[margin=25mm]{geometry}
\usepackage{fontspec}
\setmainfont{Latin Modern Roman}
\usepackage{amsmath,amssymb,mathtools}
\usepackage{unicode-math}
\setmathfont{Latin Modern Math}
\usepackage{xurl,hyperref}
\hypersetup{colorlinks=true,urlcolor=blue}
\setcounter{secnumdepth}{0}
\setlength{\parindent}{0pt}
\setlength{\parskip}{5pt}
\setlength{\emergencystretch}{3em}
\newcommand{\sref}[2]{\hyperlink{source:#1:#2}{[S#2]}}
\newcommand{\cataloguescope}[1]{\paragraph{Recorded target.}#1}
\newcommand{\R}{\mathbb R}
\newcommand{\C}{\mathbb C}
\newcommand{\N}{\mathbb N}
\newcommand{\Z}{\mathbb Z}
\newcommand{\Q}{\mathbb Q}
\newcommand{\F}{\mathbb F}
\newcommand{\eps}{\varepsilon}
\begin{document}
% UnsolvedMath ID 30001125; source catalogue code OWR-2495-002
\setcounter{section}{30001124}
\section{Critical Transition and Pattern Scales in Shear Turbulence}\label{30001125}
{\small UnsolvedMath ID 30001125 · Mathematical Physics}
\subsection{Definitions and mathematical statement}

\paragraph{Shared notation.}$\mathbb N=\{1,2,\ldots\}$; the standard number fields have their usual meanings. All additional conventions are stated locally.


\paragraph{Definitions and conventions.} The Reynolds number $\mathrm{Re}$ is the dimensionless ratio of inertial and viscous effects in the specified shear geometry. $\mathrm{Re}_c$ denotes the source's threshold where sustained turbulence becomes possible. Turbulent--laminar patterns consist of spatially varying regions of turbulent and laminar flow, with a characteristic wavelength much larger than the basic shear-layer scale.

\paragraph{Source scope and qualifications.} This threshold is not automatically the linear-instability Reynolds number: the report stresses a discontinuous subcritical transition. It asks for the mechanism and long scale of the patterns in the examined shear flows, without a universal wavelength formula.

\paragraph{Catalogue mathematical target.}
\paragraph{Catalogue wording retained for traceability.}
In shear flows, why does transition to turbulence occur near the critical Reynolds number $R_{e,c}$ so often, and why is the resulting pattern's length scale so large?

\paragraph{Saved literature context.} The supplied catalogue's assessment dated 2026-08-21 records: Shear-flow transition and the large pattern length near the critical Reynolds number remain active problems; later stability and instability results do not provide a general explanation of the observed scaling. The assessment is retained as dated provenance; it is not a new comprehensive literature or open-status certification.
\subsection{Short English statement}
Explain why many shear flows develop large turbulent--laminar patterns close to the onset of sustained turbulence.
\subsection{Sources}
\begin{itemize}
\item[\hypertarget{source:30001125:1}{S1}] ULAM.AI and the UnsolvedMath Contributors, UnsolvedMath, record 30001125 (OWR-2495-002). Catalogue curation under CC BY 4.0; definitions and English exposition added locally.
\url{https://huggingface.co/datasets/ulamai/UnsolvedMath}.
\item[\hypertarget{source:30001125:2}{S2}] Dynamics of Patterns (2008), Oberwolfach Reports, DOI 10.4171/owr/2008/57, pp. 3201--3271; official report source and mathematical context.
\url{https://doi.org/10.4171/owr/2008/57}.
\item[\hypertarget{source:30001125:3}{S3}] Official EMS Press PDF of the same Oberwolfach report; PDF pages 39--40, printed pages 3239--3240.
\url{https://ems.press/content/serial-article-files/46201}.
\end{itemize}
\label{end:30001125}
\end{document}
